\documentclass[11pt, a4paper]{article}

\usepackage[T1]{fontenc}
\usepackage[utf8]{inputenc}
\usepackage{tgpagella}
\usepackage[spanish, es-noshorthands]{babel}
\usepackage{amsmath, amssymb, bm}
\usepackage{tabularx}
\usepackage{booktabs}
\usepackage[most]{tcolorbox}
\usepackage[margin=2.5cm]{geometry}
\usepackage{ragged2e}
\usepackage{fancyhdr}

\pagestyle{fancy}
\fancyhf{}
\setlength{\headheight}{16pt}
\lhead{\textbf{UCB Sede Tarija} | Ing. Mecatrónica}
\chead{}
\rhead{IMT-121 Circuitos Electrónicos I | Prompt Microdefensa}
\lfoot{Prof: M.Sc. Ing. Bernardo Quiroga Turdera}
\rfoot{Página \thepage}
\renewcommand{\headrulewidth}{0.4pt}

\definecolor{ucbblue}{RGB}{0, 51, 102}

\begin{document}

\begin{tcolorbox}[colback=ucbblue!5!white, colframe=ucbblue, title=\textbf{Registro Selectivo de Microdefensa (Path B)}]
    \centering \textbf{Resolución de Inconsistencia de Evidencias}
\end{tcolorbox}

\textbf{Estudiante:} \rule{5cm}{0.4pt} \quad \textbf{Fecha:} \rule{4cm}{0.4pt}

\vspace{0.4cm}
\noindent\textbf{Conflicto de Evidencia (Trigger):} \\
\textit{[Ej. El estudiante determinó $R_{th}$ perfectamente en el reporte de la Exp. A, pero en el examen calculó $R_{th}$ sin apagar la fuente de tensión].} \\
\rule{\textwidth}{0.4pt} \\
\rule{\textwidth}{0.4pt}

\vspace{0.4cm}
\noindent\textbf{Pregunta 1 (Targeted Question):} \\
\textit{[Ej. Muéstreme en su diagrama del examen qué sucede físicamente con la fuente de 12V al hallar $R_{th}$ y por qué.]} \\
\rule{\textwidth}{0.4pt}

\vspace{0.2cm}
\noindent\textbf{Evidencia Esperada:} \textit{[Desactivación correcta $\rightarrow$ cortocircuito, perspectiva desde el puerto].}

\vspace{0.4cm}
\noindent\textbf{Pregunta 2 (Si es necesaria):} \\
\rule{\textwidth}{0.4pt}

\vspace{0.6cm}
\begin{tcolorbox}[colback=gray!5, colframe=gray!50, title=\textbf{Resolución (Outcome)}]
Marque la opción que corresponda tras la entrevista (Máx. 3-5 minutos):
\begin{itemize}
    \item[$\square$] \textbf{CLARIFIED:} El estudiante domina el concepto; el error en el examen fue un lapsus aislado. La evidencia del reporte se sostiene.
    \item[$\square$] \textbf{STILL\_INCONSISTENT:} El estudiante no comprende el concepto básico que afirma haber desarrollado en su reporte.
    \item[$\square$] \textbf{REQUIRES\_FOLLOW\_UP:} Posible problema de asimetría de equipo grave o sospecha de plagio en el reporte.
\end{itemize}
\end{tcolorbox}

\noindent\textbf{Notas del Instructor:} \\
\rule{\textwidth}{0.4pt} \\
\rule{\textwidth}{0.4pt}

\end{document}
